\documentclass[10pt,a4paper]{amsart}
\usepackage[margin=19mm]{geometry}
\usepackage{amsmath,amssymb,mathtools}
\usepackage[scr=rsfs,cal=euler]{mathalfa}
\usepackage{xfrac}
\usepackage[unicode,hidelinks,bookmarks=false]{hyperref}
\newcommand{\ie}{i.e.\ }
\newcommand{\cf}{cf.\ }
\newcommand{\resp}{respectively\ }
\newcommand{\Subsection}{\subsection}
\providecommand{\nobreakdash}{\nobreak\hbox{-}}
\setlength{\emergencystretch}{2em}
\multlinegap0pt
\usepackage{mathtools}
\usepackage{amssymb}
\usepackage{mdframed}
\usepackage[shortlabels]{enumitem}
\usepackage{csquotes}
\newtheorem{theorem}{Theorem}
\newtheorem{lemma}{Lemma}
\newcommand{\BF}{\mathrm{BF}}
\newcommand{\Ell}{\mathcal{E}}
\newcommand{\KS}{\mathscr{W}}
\newcommand{\Lgeom}{L_p^{\mathrm{geom}}}
\newcommand{\cBF}{\Xi}
\newcommand{\lcBF}{\widetilde{\Xi}}
\renewcommand{\le}{\leqslant}
\renewcommand{\leq}{\le}
\newcommand{\regulator}[1]{r_{#1}}
\newcommand{\rsyn}{\regulator{\syn}}
\DeclareMathOperator{\syn}{syn}
\title{On the factorisation of the \texorpdfstring{$p$}{p}-adic Rankin--Selberg \texorpdfstring{$L$}{L}-function in the supersingular case}
\author{Alessandro Arlandini, David Loeffler}
\date{Source: arXiv:2103.16380v3 (CC BY 4.0)}
\begin{document}
\maketitle

\noindent\emph{Source and licence.} On the factorisation of the \texorpdfstring{$p$}{p}-adic Rankin--Selberg \texorpdfstring{$L$}{L}-function in the supersingular case. Version arXiv:2103.16380v3 (CC BY 4.0), distributed by arXiv under the Creative Commons Attribution 4.0 International licence, http://creativecommons.org/licenses/by/4.0/, as recorded in the arXiv licence metadata for this exact version and shown on the abstract page https://arxiv.org/abs/2103.16380v3. Full licence notice: \texttt{licenses/arxiv-2103.16380-CC-BY-4.0.txt}.\par\smallskip

\noindent\textit{Editorial scope.} Counted unit: the article's theorem theorem (source0.tex lines 1757--1763). The article's complete original proof of it is delivered with it (source0.tex lines 1764--1774, one \texttt{proof} environment of 11 lines). No mathematics is written, repaired or re-derived: the statement and the proof are the article's own lines, in the article's own notation, and no retained line is substituted or re-flowed.  Retained uncounted as the source-local context the proof needs: lines 1668--1675; lines 1738--1744; lines 1748--1751.  Nothing was omitted from any retained span, so the package carries no omission record.  No retained line contains a citation, so the wrapper carries no bibliography and no bibliography entry.  No retained line contains a figure environment, an image-inclusion command or drawing code, so no image is delivered and the package declares no asset dependency.  Editorial formatting only, in addition: an AMS \texttt{amsart} preamble that restates the article's own declarations and macros used by the retained lines, namely the environments \texttt{theorem}, \texttt{lemma} on separate counters, together with the standard packages those lines need.  The retained environments print in this document as Theorem 1, Lemma 1 in order of appearance, whereas the article prints the same items with the numbering of its own full document.  The complete native source of the article as distributed in the arXiv e-print for this exact version is bundled unchanged as \texttt{sources/109730/source0.tex}.
\begin{equation}
    \label{eq:pairing-in-two}
    \begin{split}
        \langle \rsyn(\cBF^{k,k,j}),Z_f^{\syn} \rangle_{\syn, \KS_k} &= \langle \BF^{[k,k,j]}_{\syn}, Z_{f,c}^{\syn} \rangle_{\syn,\Ell^{k, k}} \\
        &+ \frac{1}{2}\langle \iota^*(i_{\mathrm{cusp},*}(\xi))_{\syn}, Z_{f,c}^{\syn} \rangle_{\syn,\Ell^{k, k}} \\
        &+ \frac{(-1)^{k+j}}{2}\langle \iota^*(i_{\mathrm{cusp},*}(\xi))_{\syn}, (\rho')^* Z_{f,c}^{\syn} \rangle_{\syn,\Ell^{k, k}}.
    \end{split}
\end{equation}

\begin{theorem}
    \label{thm:first-term}
    If $j$ is even, $0\leq j\leq k$ and $E(f,f,j+1)\neq 0$, then
    \[
        \langle \BF^{[k,k,j]}_{\syn},Z_{f,c}^{\syn}\rangle_{\syn} = (-1)^k k!\binom{k}{j}\lambda_N(f^*)\frac{2E(f)E^*(f)}{E(f,f,j+1)}\Lgeom(F,F)(k,k,j+1).
    \]
\end{theorem}

\begin{lemma}
    \label{lemma:cuspidal-is-zero}
    We have $\iota^*(i_{\mathrm{cusp},*}(\xi)) = 0$.
\end{lemma}

\begin{theorem}
    If $j$ is even, $0\leq j\leq k$ and $E(f,f,j+1)\neq 0$, then
\begin{multline*}
    \langle \rsyn(\cBF^{k,k,j}),Z_f^{\syn} \rangle_{\syn, \KS_k} = \langle \rsyn(\lcBF^{k,k,j}),Z_f^{\syn} \rangle_{\syn, \KS_k} \\
 = (-1)^k k!\binom{k}{j}\lambda_N(f^*)\frac{2E(f)E^*(f)}{E(f,f,j+1)}\Lgeom(F,F)(k,k,j+1).
\end{multline*}
\end{theorem}

\begin{proof}
    By Lemma~\ref{lemma:cuspidal-is-zero}, all the terms in equation~\eqref{eq:pairing-in-two} vanish bar the first, hence
    \[
        \langle \rsyn(\cBF^{k,k,j}),Z_f^{\syn} \rangle_{\syn, \KS_k} = \langle \BF^{[k,k,j]}_{\syn}, Z_{f,c}^{\syn} \rangle_{\syn,\Ell^{k, k}}.
    \]
    Applying Theorem~\ref{thm:first-term} shows that this computes to the $p$-adic $L$-value in the statement of the theorem. The other equality follows from the discussion after equation~\eqref{eq:pairing-in-two}. Indeed, the difference of the two pairings equals
    \[
        \frac{(-1)^{k+j}}{2}\langle \iota^*(i_{\mathrm{cusp},*}(\xi))_{\syn}, (\rho')^* Z_{f,c}^{\syn} \rangle_{\syn,\Ell^{k,k}} - \frac{1}{2}\langle \iota^*(i_{\mathrm{cusp},*}(\xi))_{\syn}, Z_{f,c}^{\syn} \rangle_{\syn,\Ell^{k,k}}
    \]
    which is again zero by the proved results.
\end{proof}

\end{document}
